% 2011 年天津卷文科第 13 题：圆中两弦 AB 与 CD 相交于 F，E 在 AB 延长线上，CE 与圆相切。
% 来源 PDF：普通高考/2011/2011天津文.pdf，第 1 页（印面 707）第 2 栏。
% 重绘依据：原图圆心约在 (0,0)，A、B 在圆上且 AB 近水平，E 在 B 的右侧圆外，
%   D 在圆上左上方，C 在圆上右下方，CD 过 F；CE 为切线，C 为切点。
\begin{tikzpicture}[scale=1.0, semithick, every node/.style={inner sep=1.5pt}]
  \coordinate (O) at (0,0);
  \draw (O) circle (2.0);
  % 弦 AB（近水平，略向右上倾斜）
  \coordinate (A) at (-160:2.0); % 左下
  \coordinate (B) at (-15:2.0);  % 右下
  \draw (A) -- (B);
  % E 在 AB 延长线上
  \coordinate (E) at (2.55,-0.28);
  \draw (B) -- (E);
  % 弦 CD 过 F（F 为 AB 与 CD 交点）
  \coordinate (F) at ($(A)!0.62!(B)$);
  \coordinate (D) at (135:2.0);
  \coordinate (C) at (-58:2.0);
  \draw (D) -- (C);
  % 切线 CE
  \draw (C) -- (E);
  % 标注
  \node[left]  at (A) {$A$};
  \node[above right] at (B) {$B$};
  \node[right] at (E) {$E$};
  \node[above left] at (D) {$D$};
  \node[below right] at (C) {$C$};
  \node[above] at (F) {$F$};
\end{tikzpicture}
